% Drawing 141/01-9 -- folder 141, batch 01, archive page 10 (PDF page 11).
% Transcription: en bas de la page, une carte en forme de pentagone : cinq sommets cerclés reliés en cycle, une arête pendante issue de chacun (celle du haut finissant en fourche), un point cerclé au centre ; à droite, un astérisque.
% Kind: pentagonal map with ringed vertices and pendant edges; an asterisk. Method 2 (schematic redraw from the facsimile). Status: ok.
% Notes: The pentagon between its five ringed vertices (points with an open ring round each; the left one larger), the five sides, one pendant edge from each vertex (to the left with an end point; down to the left with an end point; down to the right; up to the right with an end point; and at the top a short edge to a small vertex that forks into two prongs and sends another edge to the left, to an end point above the left vertex), and a ringed point in the middle of the face. Below that end point, after a small gap, a short stroke goes down to the ring of the left vertex. At the right, apart, an asterisk of three strokes through one point. Left out: the frame of the boxed formula and its text, which run into the crop at the top right. Overlaid 2026-11-12 (Opus 5.5) on the 400 dpi colour crop and corrected: the right vertex was drawn 7.5 units too high (at the level of the end of its pendant edge), which bent the top and right sides; the lower right vertex, the end of the right pendant edge and the rings are moved onto the page; the short stroke under the upper left end point is drawn solid with its gap, as on the page, not dashed. Redrawn 2026-11-12 (issue 128): to check again.
% Coordinates: 1 unit = 1 mm = 1/100 of the crop width (bbox in page fractions: [0.22, 0.7, 0.58, 0.87]).
% Redrawn by Claude (claude-fable-5-2), issue #128, 2026-11-12; not checked by a human; not compiled (no TeX in the session).
% Overlaid by Claude (claude-opus-5-5) 2026-11-12 (issue #128): rasterised from its own coordinates over the 400 dpi colour facsimile and corrected from it. Not compiled. Status: ok.
% Rights: reproduces part of the work -- behind RENDER_TRANSCRIPTS, never published by itself.
\documentclass[tikz,border=3pt]{standalone}
\usepackage{amsmath,amssymb}
\usetikzlibrary{patterns,arrows.meta}
\tikzset{pencil/.style={black!45}, struck/.style={black!35}, mpt/.style={circle,draw,inner sep=1.3pt}, pt/.style={circle,fill,inner sep=0.9pt}, opt/.style={circle,draw,fill=white,inner sep=1.1pt}}
\begin{document}
\begin{tikzpicture}[x=1mm,y=1mm,line cap=round,line join=round,font=\scriptsize]
\draw (37,-22) -- (61,-34.5) -- (56.5,-55.5) -- (30.5,-49) -- (24.5,-32.5) -- cycle;
\node[pt] at (37,-22) {};
\draw (37,-21.7) circle[radius=2.9];
\node[pt] at (24.5,-32.5) {};
\draw (24.5,-32.5) circle[radius=4.4];
\node[pt] at (30.5,-49) {};
\draw (31,-49) circle[radius=3];
\node[pt] at (56.5,-55.5) {};
\draw (57.3,-56) circle[radius=2.8];
\node[pt] at (61,-34.5) {};
\draw (61.5,-34.3) circle[radius=3.4];
\node[pt] at (42.5,-38.5) {};
\draw (42.5,-38.5) circle[radius=2.8];
\draw (24.5,-32.5) -- (5,-30);
\node[pt] at (5,-30) {};
\draw (30.5,-49) -- (26,-59.5);
\node[pt] at (26,-59.5) {};
\draw (56.5,-55.5) -- (62,-62);
\draw (61,-34.5) -- (74,-27);
\node[pt] at (74,-27) {};
\draw (37,-22) -- (31.3,-16.3);
\node[pt] at (31.3,-16.3) {};
\draw plot[smooth, tension=0.6] coordinates {(31.3,-16.3) (30.5,-12) (29,-8.5)};
\draw (31.3,-16.3) -- (36,-9.5);
\draw (31.3,-16.3) -- (22.5,-20);
\node[pt] at (22.5,-20) {};
\draw (23.2,-22.5) -- (24.4,-28);
\draw (84,-40) -- (94,-42);
\draw (87.3,-37) -- (88.6,-46);
\draw (82.5,-46) -- (93,-37.5);
\end{tikzpicture}
\end{document}
